
Cross-repository metadata heterogeneity limits automated dataset comparison across HuggingFace, OpenML, and UCI repositories. This study tested whether lifecycle-stage functional categories from the Datasheets for Datasets framework---specifically the distinction between General Information and Responsible AI documentation---manifest as unsupervised cluster structure in semantic embeddings. On 300 metadata fields collected from three repositories, linear probes achieved 79.6\% accuracy in distinguishing lifecycle categories, demonstrating that distributional signals exist. However, K-means clustering achieved normalized mutual information (NMI) of 0.023, failing to exceed baseline methods and falling 96\% short of the threshold for meaningful recovery. This supervised-unsupervised performance gap indicates that lifecycle categories are encoded in embeddings as detectable signals but do not form natural clusters. The failure is attributed to severe class imbalance (8.3\% Responsible AI fields, 11:1 ratio) and cross-repository heterogeneity. Repository-stratified analysis revealed 30-fold NMI variance (UCI: 0.394, HuggingFace: 0.013), correlating with differences in class balance across repositories. These results indicate that label-free cross-repository lifecycle detection via unsupervised clustering is not feasible under current metadata distributions, but supervised or semi-supervised approaches may enable practical deployment.

